\subsection{Interprétation mécanique}
Le pneu ne porte que sur une fraction de la surface : 30 à 85\,\% selon la texture. La pression
réelle est donc concentrée sur les sommets, à 3 à 9 fois la pression nominale. Cette surpression
se diffuse dans le massif sur une profondeur de l'ordre de la taille des zones de contact, soit
quelques millimètres, et la longueur d'onde de la texture fixe la profondeur d'influence.
Au-delà, seule compte la résultante, identique au cas lisse (principe de Saint-Venant). C'est
pourquoi le TFE, qui a montré une sensibilité croissante à l'empreinte quand on se rapproche de
la surface, trouve ici un prolongement naturel : à l'échelle du pneu, la forme de l'empreinte
agit sur les premiers centimètres ; à l'échelle de la texture, les sommets agissent sur les
premiers millimètres.

La viscoélasticité joue un rôle non trivial. Les sommets de texture sont chargés dès l'arrivée
du pneu et pendant tout son passage, et la roue avant laisse une déformation différée au même
endroit de la chaussée avant le passage de la roue arrière. La méthode héréditaire en tient
compte exactement ; l'ancienne formule multiplicative sous-estimait l'effet des textures aléatoires d'environ 9\,\%.

\subsection{Limites}
Les trois limites structurelles du niveau 1 (H3, H6, H7) sont levées par le niveau 2 ; celles qui
restent sont les suivantes.
\begin{enumerate}[nosep]
\item \textbf{Caoutchouc.} Le module de la bande de roulement du pneu d'A340 (et sa
      dépendance à la fréquence et à la température) reste le paramètre le plus influent et le
      moins connu. Données à obtenir auprès du manufacturier ou dans la littérature (Michelin,
      Airbus).
\item \textbf{Textures synthétiques.} Elles sont réalistes dans leurs statistiques mais ne
      remplacent pas des relevés réels. Le code accepte directement un relevé de hauteurs (laser
      3D) à la place de la surface générée.
\item \textbf{Modèles locaux 2D.} L'entaille et la microstructure sont traitées en coupe
      (déformation plane généralisée), avec des granulats circulaires et un mortier élastique au
      point de fonctionnement. Une cellule 3D (granulats anguleux, mortier viscoélastique) est
      l'étape suivante naturelle ; le solveur EF s'étend au 3D sans changement de principe.
\item \textbf{Frottement} en glissement total (loi d'Amontons locale) : majorant pour la
      texture, et représentatif d'un freinage appuyé. Le virage (effort transversal) n'a pas été
      calculé ; \cs{} le traite déjà.
\item \textbf{Un seul cas de chargement} : vitesse et température du TFE. À vitesse plus élevée
      ou à température plus basse, l'enrobé est plus raide : la déformation diminue, mais pas
      nécessairement la contrainte, et le contraste granulat/mortier diminue.
\item \textbf{Critère d'endommagement.} Le rapport compare des déformations. Relier la
      déformation locale dans le mortier à l'amorçage (fissuration par le haut, arrachement)
      demande un critère calé sur essais (mortier bitumineux, essais de fatigue de mortier).
\end{enumerate}

\subsection{Suite proposée}
\begin{enumerate}[nosep]
\item \textbf{Relevés de texture réels} sur piste (STAC, profilomètre laser), calage de la DSP et
      comparaison avec les surfaces synthétiques.
\item \textbf{Module de la bande de roulement} : données pneumatiques, ou identification inverse
      par des mesures de pression de contact sur surface texturée (capteurs en film).
\item \textbf{Cellule mésoscopique 3D} à partir d'une tomographie d'enrobé réel, mortier
      viscoélastique ; comparaison avec la cellule 2D pour chiffrer l'effet de coupe.
\item \textbf{Essais sur mortier bitumineux} (module complexe et fatigue) pour transformer
      $K_m$ en critère d'amorçage superficiel.
\item \textbf{Indicateur mécanique de texture} : relier l'amplification à MPD, $m_2$ et $S_{sk}$
      sur une famille de surfaces, et proposer un indicateur utilisable en ingénierie.
\item \textbf{Confrontation aux dégradations observées} : fissuration par le haut et arrachement
      aux bords de rainures (Toulouse-Blagnac).
\end{enumerate}
